\chapter{Pràctica}
\label{cap:practica}

\section{Lògica}

\begin{exercici}
Reformuleu els enunciats següents fent ús de variables i constants:
(a) El producte de dos nombres parells és parell; (b) Donat un nombre real
no negatiu, busqueu dos nombres reals la diferència dels seus quadrats no
sigui més gran que el nombre donat.
\end{exercici}

\begin{solucio}
(a) Si simbolitzem per $P$ el predicat \enquote{ser
parell}, aleshores podem escriure:%
\[
\forall x,y\left(  P\left(  x\right)  \wedge P\left(  y\right)
\longrightarrow P\left(  x\cdot y\right)  \right)  \text{.}%
\]
Observa que també podem escriure:%
\[
\forall x\left(  P\left(  x\right)  \longleftrightarrow\left(  \exists
k\right)  \left(  k\in\mathbb{Z}\wedge x=2k\right)  \right)  \text{.}%
\]


(b) Suposem donat un nombre real $\delta\geq0$, o sigui no negatiu, aleshores
podem escriure:%
\[
\exists x,y\left(  x,y\in\mathbb{R}\wedge x^{2}-y^{2}\leq\delta\right)
\text{.}%
\]

\end{solucio}

\begin{exercici}
Quines de les expressions següents són proposicions i quines
predicats? (a) $4<3$; (b) $y>7$; (c) $x+y=z$.
\end{exercici}

\begin{solucio}
(a) És una proposició falsa, perquè és un enunciat del qual podem determinar el valor de veritat; (b)
és un predicat perquè és un enunciat que conté una variable
$y$ i si li assignem un valor es té una proposició; (c) És un
predicat de tres variables $x,y$ i $z$. Assignant valors a les tres variables
s'obté una proposició.
\end{solucio}

\begin{exercici}
Quins dels següents enunciats són predicats i quins no ho són?
Justifica la teva resposta. (a) $x$ és un divisor de 210; (b) El producte
dels nombres $x$ i $y$; (c) La suma de dos nombres és menor que $1$.
\end{exercici}

\begin{solucio}
(a) És un predicat d'una variable $x$ que és vertader quan, per
exemple, $x$ pren el valor $7$, i fals quan, per exemple, és $8$; (b) No
és predicat perquè quan assignem valors a les variables $x$ i $y$ no
s'obté una proposició. De fet, l'enunciat és l'expressió
algebraica $x\cdot y$; (c) És un predicat que conté dues variables.
Podem expressar-lo com $x+y<1$. És clar que quan donem valors a $x$ i $y$
s'obté una proposició.
\end{solucio}

\begin{exercici}
Considereu els enunciats següents: $p=$ \enquote{Estic
content}, $q=$ \enquote{Estic veient una
pel·lícula}\ i $r=$ \enquote{Estic menjant espagueti}. Expresseu en paraules els enunciats
següents: (a) $\left(  q\vee r\right)  \longrightarrow p$; (b) $\left(
p\wedge\lnot q\right)  \longleftrightarrow\left(  q\vee r\right)  $.
\end{exercici}

\begin{solucio}
(a) Estic content quan veig una pel·lícula o mengo
espagueti; (b) Veig una pel·lícula o mengo espagueti
només si estic content sense veure una pel·lícula.
\end{solucio}

\begin{exercici}
Analitza les formes lògiques dels enunciats següents: (a) El joc es
cancel·larà si plou o neva; (b) Tenir almenys deu
persones és una condició necessària i suficient per a la
conferència que s'està impartint; (c) Si en Miquel va anar a la
botiga, llavors tenim ous a casa, si no no en tenim.
\end{exercici}

\begin{solucio}
(a) Si $C=$\enquote{El joc serà cancel·
lat}, $P=$\enquote{Plou}\ i
$N=$\enquote{Està nevant}. Aleshores l'enunciat \enquote{El joc es cancel·larà si plou o neva} és la proposició $(P\vee N)\longrightarrow C$.

(b) Si $P=$\enquote{Hi ha almenys deu persones}\ i
$C=$\enquote{La conferència es donarà}.
Aleshores l'enunciat \enquote{Tenir almenys deu persones és una
condició necessària i suficient per a la conferència que
s'està impartint}\ és la proposició
$P\longleftrightarrow C$.

(c) Si $B=$\enquote{Miquel va anar a la botiga} i
$C=$\enquote{Hi ha ous a casa}. Aleshores, l'enunciat
\enquote{Si en Miquel va anar a la botiga, llavors tenim ous a casa,
si no no en tenim}\ és la proposició $(B\rightarrow
C)\wedge(\lnot B\rightarrow\lnot C)$. Això és equivalent a
$B\longleftrightarrow C$ perquè per la llei del contrarecíproc
s'obté $(B\rightarrow C)\wedge(C\rightarrow B)$, què és alhora
equivalent al que hem dit.
\end{solucio}

\begin{exercici}
Considereu els enunciats següents: $p=$ \enquote{L'Eduard té
els cabells vermells}, $q=$ \enquote{L'Eduard té
un nas gran}\ i $r=$ \enquote{A l'Eduard li agrada
menjar crispetes}. Tradueix els enunciats següents a
símbols: (a) \enquote{L'Eduard té els cabells vermells i no
té un nas gran}; (b) \enquote{No és el cas
que l'Eduard tingui un nas gran o li agradi menjar crispetes}.
\end{exercici}

\begin{solucio}
(a) $p\wedge\lnot q$; (b) $\lnot(q\vee r)$.
\end{solucio}

\begin{exercici}
Construeix les taules de veritat de les següents expressions lògiques:
(a) $(p\wedge q)\vee\lnot p$; (b) $\lnot p\longrightarrow\lnot(q\vee r)$.
\end{exercici}

\begin{solucio}
(a)%
\[%
\begin{tabular}
[c]{ll|l|ll}%
$p$ & $q$ & $\lnot p$ & $p\wedge q$ & $(p\wedge q)\vee\lnot p$\\\hline
\multicolumn{1}{c}{V} & \multicolumn{1}{c|}{V} & F & \multicolumn{1}{|c}{V} &
V\\
\multicolumn{1}{c}{V} & \multicolumn{1}{c|}{F} & F & \multicolumn{1}{|c}{F} &
F\\
\multicolumn{1}{c}{F} & \multicolumn{1}{c|}{V} & V & \multicolumn{1}{|c}{F} &
V\\
\multicolumn{1}{c}{F} & \multicolumn{1}{c|}{F} & V & \multicolumn{1}{|c}{F} &
V
\end{tabular}
\]


(b)%
\[%
\begin{tabular}{ccc|c|c|c}
$p$ & $q$ & $r$ & $\lnot p$ & $\lnot(q\vee r)$ & $\lnot p\longrightarrow\lnot(q\vee r)$\\\hline
V & V & V & F & F & V\\
V & V & F & F & F & V\\
V & F & V & F & F & V\\
V & F & F & F & V & V\\
F & V & V & V & F & F\\
F & V & F & V & F & F\\
F & F & V & V & F & F\\
F & F & F & V & V & V
\end{tabular}
\]

\end{solucio}

\begin{exercici}
Considerem els predicats $P(x,y)=$ \enquote{$x+y=4$} i $Q(x,y)=$ \enquote{$x<y$}. Troba els
valors de veritat de les següents expressions lògiques: (a) $\lnot
P(x,y)\vee Q(x,y)$ i (b) $\lnot(P(x,y)\longleftrightarrow Q(x,y))$ quan usem
els valors següents (1) $x=1,y=3$; (2) $x=2,y=1$.
\end{exercici}

\begin{solucio}
Per a l'expressió $\lnot P(x,y)\vee Q(x,y)$, tenim:
(1) amb $x=1$ i $y=3$, és vertadera, ja que $Q(1,3)$ és vertadera;
(2) amb $x=2$ i $y=1$, també és vertadera, ja que $P(2,1)$ és falsa.

Per a l'expressió $\lnot(P(x,y)\longleftrightarrow Q(x,y))$, tenim:
(1) amb $x=1$ i $y=3$, és falsa, ja que $P(1,3)$ i $Q(1,3)$ són
ambdues vertaderes; (2) amb $x=2$ i $y=1$, també és falsa, ja que
$P(2,1)$ i $Q(2,1)$ són ambdues falses.
\end{solucio}

\begin{exercici}
Proveu les implicacions següents: (a) $(A\longrightarrow B)\wedge\lnot
B\Longrightarrow\lnot A$; (b) $\left(  A\longrightarrow B\right)
\wedge\left(  B\longrightarrow C\right)  \Longrightarrow A\longrightarrow C$.
\end{exercici}

\begin{solucio}
(a) Per veure que $(A\longrightarrow B)\wedge\lnot B\Longrightarrow\lnot A$
hem de comprovar que $\left(  (A\longrightarrow B)\wedge\lnot B\right)  $
$\longrightarrow\lnot A$ és tautologia. Construïm la taula de veritat
corresponent:%
\[%
\begin{tabular}
[c]{ccccccc}%
$A$ & $B$ & $\lnot A$ & $\lnot B$ & $A$ $\longrightarrow$ $B$ &
$(A\longrightarrow B)\wedge\lnot B$ & $\left(  (A\longrightarrow B)\wedge\lnot
B\right)  $ $\longrightarrow\lnot A$\\\hline
V & V & F & F & V & F & V\\
V & F & F & V & F & F & V\\
F & V & V & F & V & F & V\\
F & F & V & V & V & V & V
\end{tabular}
\]
i s'obté una tautologia doncs veiem que l'última columna només
té el valor V.

(b) Per veure que $\left(  A\longrightarrow B\right)  \wedge\left(
B\longrightarrow C\right)  \Longrightarrow A\longrightarrow C$ hem de
comprovar que $\left(  \left(  A\longrightarrow B\right)  \wedge\left(
B\longrightarrow C\right)  \right)  $ $\longrightarrow\left(  A\longrightarrow
C\right)  $ és tautologia. Construïm la taula de veritat corresponent
prenent $\alpha=A\longrightarrow B$, $\beta=B\longrightarrow C $ i
$\gamma=A\longrightarrow C$:%
\[%
\begin{tabular}
[c]{cccccccc}%
$A$ & $B$ & $C$ & $\alpha$ & $\beta$ & $\gamma$ & $\alpha\wedge\beta$ &
$\left(  \alpha\wedge\beta\right)  \longrightarrow\gamma$\\\hline
V & V & V & V & V & V & V & V\\
V & V & F & V & F & F & F & V\\
V & F & V & F & V & V & F & V\\
V & F & F & F & V & F & F & V\\
F & V & V & V & V & V & V & V\\
F & V & F & V & F & V & F & V\\
F & F & V & V & V & V & V & V\\
F & F & F & V & V & V & V & V
\end{tabular}
\]
i surt que és tautologia com era d'esperar.
\end{solucio}

\begin{exercici}
Proveu les equivalències lògiques següents: (a) $A\longrightarrow
B\Longleftrightarrow\lnot A\vee B$; (b) $A\longrightarrow\left(  B\wedge
C\right)  \Longleftrightarrow\left(  A\longrightarrow B\right)  \wedge\left(
A\longrightarrow C\right)  $.
\end{exercici}

\begin{solucio}
(a) Per veure que $A\longrightarrow B\Longleftrightarrow\lnot A\vee B$ hem de
comprovar que $(A\longrightarrow B)\longleftrightarrow\left(  \lnot A\vee
B\right)  $ és tautologia. Construïm la taula de veritat
corresponent:%
\[%
\begin{tabular}
[c]{cccccc}%
$A$ & $B$ & $\lnot A$ & $A$ $\longrightarrow$ $B$ & $\lnot A\vee B$ &
$(A\longrightarrow B)\longleftrightarrow\left(  \lnot A\vee B\right)
$\\\hline
V & V & F & V & V & V\\
V & F & F & F & F & V\\
F & V & V & V & V & V\\
F & F & V & V & V & V
\end{tabular}
\]
i s'obté una tautologia doncs veiem que l'última columna només
té el valor V.

(b) Per veure que $A\longrightarrow\left(  B\wedge C\right)
\Longleftrightarrow\left(  A\longrightarrow B\right)  \wedge\left(
A\longrightarrow C\right)  $ hem de comprovar que $(A\longrightarrow\left(
B\wedge C\right)  )\longleftrightarrow\left(  \left(  A\longrightarrow
B\right)  \wedge\left(  A\longrightarrow C\right)  \right)  $ és
tautologia. Construïm la taula de veritat corresponent, prenent
$\alpha=A\longrightarrow B$, $\beta=A\longrightarrow C$ i $\gamma
=A\longrightarrow\left(  B\wedge C\right)  $:
\[%
\begin{tabular}
[c]{ccccccccc}%
$A$ & $B$ & $C$ & $\alpha$ & $B\wedge C$ & $\beta$ & $\gamma$ & $\delta
=\alpha\wedge\beta$ & $\gamma\longleftrightarrow\delta$\\\hline
V & V & V & V & V & V & V & V & V\\
V & V & F & V & F & F & F & F & V\\
V & F & V & F & F & V & F & F & V\\
V & F & F & F & F & F & F & F & V\\
F & V & V & V & V & V & V & V & V\\
F & V & F & V & F & V & V & V & V\\
F & F & V & V & F & V & V & V & V\\
F & F & F & V & F & V & V & V & V
\end{tabular}
\]
i surt que és tautologia com era d'esperar.
\end{solucio}

\begin{exercici}
Simplifiqueu les afirmacions següents. Podeu fer ús de les
equivalències de l'exercici anterior a més de les equivalències
comentades a teoria i pràctica: (a) $X\longrightarrow\left(  X\wedge
Y\right)  $; (b) $\lnot\left(  X\vee Y\right)  \wedge Y$.
\end{exercici}

\begin{solucio}
(a) Aplicant la primera equivalència de l'apartat anterior es té:
$X\longrightarrow\left(  X\wedge Y\right)  \Longleftrightarrow\lnot
X\vee\left(  X\wedge Y\right)  $. Aplicant ara la llei distributiva de $\vee$
respecte de $\wedge$ s'obté%
\[
\lnot X\vee\left(  X\wedge Y\right)  \Longleftrightarrow\left(  \lnot X\vee
X\right)  \wedge\left(  \lnot X\vee Y\right)
\]
Ara bé, $\lnot X\vee X$ és tautologia. Per tant, $\left(  \lnot X\vee
X\right)  \wedge\left(  \lnot X\vee Y\right)  \Longleftrightarrow\lnot X\vee Y
$. Com a conseqüència, obtenim que $X\longrightarrow\left(  X\wedge Y\right)
\Longleftrightarrow X\longrightarrow Y$, doncs $\lnot X\vee
Y\Longleftrightarrow X\longrightarrow Y$.

(b) Aplicant les lleis de De Morgan es té:%
\[
\lnot\left(  X\vee Y\right)  \wedge Y\Longleftrightarrow\left(  \lnot
X\wedge\lnot Y\right)  \vee Y.
\]
Ara, aplicant la llei distributiva de $\vee$ respecte de $\wedge$, s'obté
\[
\lnot\left(  X\vee Y\right)  \wedge Y\Longleftrightarrow\left(  \lnot X\vee
Y\right)  \wedge\left(  \lnot Y\vee Y\right)  .
\]
Finalment, com $\lnot Y\vee Y$ és tautologia i $\lnot X\vee
Y\Longleftrightarrow X\longrightarrow Y$, es té%
\[
\lnot\left(  X\vee Y\right)  \wedge Y\Longleftrightarrow X\longrightarrow
Y\text{.}%
\]

\end{solucio}

\begin{exercici}
Escriu la negació dels enunciats següents: (a) $3>5$ o $7\leq8$; (b)
Si $x=3$, llavors $x^{2}=5$; (c)\ $a-b=c$ si i només si $a=b+c$.
\end{exercici}

\begin{solucio}
(a) Aplicant la llei de de Morgan, $3\leq5$ o $7>8$.

(b) L'enunciat és de la forma $A\longrightarrow B$, aleshores la seva
negació és $\lnot\left(  A\longrightarrow B\right)  \,$. Ara bé,
sabem que $A\longrightarrow B\Longleftrightarrow\lnot A\vee B$ i aplicant de
nou la llei de De Morgan es té $\lnot\left(  A\longrightarrow B\right)
\Longleftrightarrow A\wedge\lnot B$. Per tant, la negació de l'enunciat
és $x=3$ i $x^{2}\neq5$.

(c) Observem primer que $a-b=c$ és equivalent a $a=b+c$. Per tant,
l'enunciat és tautològic i de la forma $A\longleftrightarrow A$,
aleshores la seva negació és contradicció: $a-b=c$ i $a\neq b+c$.
\end{solucio}

\section{Raonament lògic}

\begin{exercici}
Per a cadascun dels raonaments següents, si és vàlid, doneu una
deducció, i si no és vàlid, mostreu el perquè.

\begin{enumerate}
\item Si el fre falla i el camí està gelat, llavors el cotxe no
pararà. Si el cotxe es va revisar, no fallaran els frens. Però el
cotxe no es va revisar. Per tant, el cotxe no pararà.

\item Si el punt d'una recta representa un enter, llavors el número es pot
definir com un decimal infinit o per un parell de decimals infinits. O el
número es pot definir per un decimal finit o el número pot ser definit
o bé per un decimal infinit o per un parell de decimals infinits. El
número no pot ser definit per un decimal finit. Per tant, el punt de la
recta representa un enter.

\item Si el rellotge està avançat, llavors en Joan va arribar abans de
les deu i va veure sortir el cotxe de l'Andreu. Si l'Andreu no diu la veritat,
llavors en Joan no va veure sortir el cotxe de l'Andreu. L'Andreu no diu la
veritat o estava en l'edifici en el moment del crim. El rellotge està
avançat. Per tant, l'Andreu estava en l'edifici en el moment del crim.
\end{enumerate}
\end{exercici}

\begin{solucio}
(1) Les proposicions simples d'aquest raonament són: $p=$\enquote{els frens fallen}, $q=$\enquote{el camí
està gelat}, $r=$\enquote{el cotxe parará} i
$s=$\enquote{el cotxe es va revisar}. Expressem ara
el raonament simbòlicament:%
\[%
\begin{tabular}
[c]{ll}%
$P_{1}:$ & $\left(  p\wedge q\right)  \longrightarrow\lnot r$\\
$P_{2}:$ & $s\longrightarrow\lnot p$\\
$P_{3}:$ & $\lnot s$\\\hline
$C:$ & $\lnot r$%
\end{tabular}
\]


L'argument és fals. Prenem com interpretació: $%
\begin{tabular}
[c]{cccc}%
$p$ & $q$ & $r$ & $s$\\\hline
$V$ & $F$ & $V$ & $F$%
\end{tabular}
$, aleshores les premisses són vertaderes
\begin{tabular}
[c]{ccc}%
$\left(  p\wedge q\right)  \longrightarrow\lnot r$ & $s\longrightarrow\lnot p
$ & $\lnot s$\\\hline
$V$ & $V$ & $V$%
\end{tabular}
i la conclusió, falsa
\begin{tabular}
[c]{c}%
$\lnot r$\\\hline
$F$%
\end{tabular}
.

(2) Les proposicions simples d'aquest raonament són: $p=$\enquote{el punt d'una recta representa un enter}, $q=$%
\enquote{el número es pot definir com un decimal
infinit}, $r=$\enquote{el número es pot definir
per un parell de decimals infinits} i $s=$\enquote{el número es
pot definir per un decimal finit}. Expressem ara el raonament
simbòlicament:%
\[%
\begin{tabular}
[c]{ll}%
$P_{1}:$ & $p\longrightarrow\left(  q\vee r\right)  $\\
$P_{2}:$ & $\lnot\left(  s\longleftrightarrow\left(  q\vee r\right)  \right)
$\\
$P_{3}:$ & $\lnot s$\\\hline
$C:$ & $p$%
\end{tabular}
\]


L'argument és fals. Prenem com interpretació: $%
\begin{tabular}
[c]{cccc}%
$p$ & $q$ & $r$ & $s$\\\hline
$F$ & $V$ & $V$ & $F$%
\end{tabular}
$, aleshores les premisses són vertaderes
\begin{tabular}
[c]{ccc}%
$p\longrightarrow\left(  q\vee r\right)  $ & $\lnot\left(
s\longleftrightarrow\left(  q\vee r\right)  \right)  $ & $\lnot s$\\\hline
$V$ & $V$ & $V$%
\end{tabular}
i la conclusió, falsa
\begin{tabular}
[c]{c}%
$p$\\\hline
$F$%
\end{tabular}
.

(3) Les proposicions simples d'aquest raonament són: $p=$\enquote{el rellotge està avançat}, $q=$\enquote{Joan
va arribar abans de les deu}, $r=$\enquote{Joan va
veure sortir el cotxe de l'Andreu}, $s=$\enquote{l'Andreu diu la
veritat} i $t=$\enquote{Andreu estava en l'edifici
en el moment del crim}. Expressem ara el raonament
simbòlicament:%
\[%
\begin{tabular}
[c]{ll}%
$P_{1}:$ & $p\longrightarrow\left(  q\wedge r\right)  $\\
$P_{2}:$ & $\lnot s\longrightarrow\lnot r$\\
$P_{3}:$ & $\lnot s\vee t$\\
$P_{4}:$ & $p$\\\hline
$C:$ & $t$%
\end{tabular}
\]
Ara provarem la validesa d'aquest argument utilitzant les regles
d'inferència:%
\[%
\begin{tabular}
[c]{lll}%
1. & $p\longrightarrow\left(  q\wedge r\right)  $ & $P_{1}$\\
2. & $\lnot s\longrightarrow\lnot r$ & $P_{2}$\\
3. & $\lnot s\vee t$ & $P_{3}$\\
4. & $p$ & $P_{4}$\\
5. & $q\wedge r$ & MP(1,4)\\
6. & $r$ & EC 5\\
7. & $\lnot\lnot r$ & DN 6\\
8. & $\lnot\lnot s$ & MT(2,7)\\\hline
9. & $t$ & ED(3,8)
\end{tabular}
\
\]

\end{solucio}

\begin{exercici}
Per a cadascun dels arguments següents, si és vàlid, doneu una
deducció, i si no és vàlid, mostreu el perquè.

\begin{enumerate}
\item $%
\begin{tabular}
[c]{ll}%
$P_{1}:$ & $p\longrightarrow q$\\
$P_{2}:$ & $\lnot r\longrightarrow\left(  t\longrightarrow s\right)  $\\
$P_{3}:$ & $r\vee\left(  p\vee t\right)  $\\
$P_{4}:$ & $\lnot r$\\\hline
$C:$ & $q\vee s$%
\end{tabular}
\ \ \ \ \ \ $

\item $%
\begin{tabular}
[c]{ll}%
$P_{1}:$ & $\lnot p\longrightarrow\left(  q\longrightarrow\lnot r\right)  $\\
$P_{2}:$ & $r\longrightarrow\lnot p$\\
$P_{3}:$ & $\left(  \lnot s\vee p\right)  \longrightarrow\lnot\lnot r$\\
$P_{4}:$ & $\lnot s$\\\hline
$C:$ & $\lnot q$%
\end{tabular}
$
\end{enumerate}
\end{exercici}

\begin{solucio}
(1) L'argument és vàlid i una deducció és:
\[%
\begin{tabular}
[c]{lll}%
1. & $p\longrightarrow q$ & $P_{1}$\\
2. & $\lnot r\longrightarrow\left(  t\longrightarrow s\right)  $ & $P_{2}$\\
3. & $r\vee\left(  p\vee t\right)  $ & $P_{3}$\\
4. & $\lnot r$ & $P_{4}$\\
5. & $p\vee t$ & ED(3,4)\\
6. & $t\longrightarrow s$ & MP(2,4)\\\hline
7. & $q\vee s$ & DL(1,5,6)
\end{tabular}
\]


(2) L'argument també és vàlid i una deducció és:%
\[%
\begin{tabular}
[c]{lll}%
1. & $\lnot p\longrightarrow\left(  q\longrightarrow\lnot r\right)  $ &
$P_{1}$\\
2. & $r\longrightarrow\lnot p$ & $P_{2}$\\
3. & $\left(  \lnot s\vee p\right)  \longrightarrow\lnot\lnot r$ & $P_{3}$\\
4. & $\lnot s$ & $P_{4}$\\
5. & $\lnot s\vee p$ & ID 4\\
6. & $\lnot\lnot r$ & MP(3,5)\\
7. & $r$ & DN 6\\
8. & $\lnot p$ & MP(2,7)\\
9. & $q\longrightarrow\lnot r$ & MP(1,8)\\\hline
10. & $\lnot q$ & MT(6,9)
\end{tabular}
\]

\end{solucio}

\begin{exercici}
Escriviu una deducció per el següent raonament. Indiqueu si les
premisses són consistents o inconsistents:\enquote{Els ordinadors
són útils i divertits i els ordinadors consumeixen molt de temps. Si
els ordinadors són difícils d'utilitzar, no són divertits. Si els
ordinadors no estan ben dissenyats, llavors són difícils d'utilitzar.
Per tant, els ordinadors estan ben dissenyats}.
\end{exercici}

\begin{solucio}
(1) Les proposicions simples d'aquest raonament són: $p=$\enquote{els ordinadors són útils i divertits}, $q=$%
\enquote{els ordinadors consumeixen molt de temps},
$r=$\enquote{els ordinadors són difícils d'utilitzar} i
$s=$\enquote{els ordinadors estan ben dissenyats}.
Expressem ara el raonament simbòlicament:%
\[%
\begin{tabular}
[c]{ll}%
$P_{1}:$ & $p\wedge q$\\
$P_{2}:$ & $r\longrightarrow\lnot p$\\
$P_{3}:$ & $\lnot s\longrightarrow r$\\\hline
$C:$ & $s$%
\end{tabular}
\]


Ara provarem la validesa d'aquest argument utilitzant les regles
d'inferència:%
\[%
\begin{tabular}
[c]{lll}%
1. & $p\wedge q$ & $P_{1}$\\
2. & $r\longrightarrow\lnot p$ & $P_{2}$\\
3. & $\lnot s\longrightarrow r$ & $P_{3}$\\
5. & $p$ & EC(1)\\
6. & $\lnot\lnot p$ & DN 5\\
7. & $\lnot r$ & MT(2,6)\\
8. & $\lnot\lnot s$ & MT(3,7)\\\hline
9. & $s$ & DN 8
\end{tabular}
\]


Podem provar la consistència de les premisses comprovant que hi ha una
interpretació que faci totes les premisses vertaderes:%
\[%
\begin{tabular}
[c]{lll|llll|llll}%
$p$ & $\wedge$ & $q$ & $r$ & $\longrightarrow$ & $\lnot$ & $p$ & $\lnot$ & $s$
& $\longrightarrow$ & $r$\\
& V &  &  & V &  &  &  &  & V & \\
V & V & V & F & V & F & V & F & V & V & F
\end{tabular}
\
\]

\end{solucio}

\begin{exercici}
És una fal·làcia el raonament següent:
\enquote{Si el poble elegeix un alcalde, s'augmentaran els impostos.
Si no s'augmenten els impostos al poble, llavors no es construirà un
estadi nou. El poble no tria un alcalde. Per tant, no es construirà un
estadi nou}.
\end{exercici}

\begin{solucio}
Si formalitzem l'argument prenent: $p=$\enquote{el poble elegeix un
alcalde}, $q=$\enquote{augmentarà els
impostos}, i $r=$\enquote{es construirà un
estadi nou}. Aleshores, es té%
\[%
\begin{tabular}
[c]{ll}%
$P_{1}:$ & $p\longrightarrow q$\\
$P_{2}:$ & $\lnot q\longrightarrow\lnot r$\\
$P_{3}:$ & $\lnot p$\\\hline
$C:$ & $\lnot r$%
\end{tabular}
\]
Es clar que es tracta d'una fal·làcia perquè es
pensa que pel fet de no triar un alcalde no s'augmentaran els impostos i com a
conseqüència no es construirà un estadi nou. Això és fals com
es pot veure a continuació:%
\[%
\begin{tabular}
[c]{lll|lllll|ll|ll}%
$p$ & $\longrightarrow$ & $q$ & $\lnot$ & $q$ & $\longrightarrow$ & $\lnot$ &
$r$ & $\lnot$ & $p$ & $\lnot$ & $r$\\
& V &  &  &  & V! &  &  & V &  & F & \\
F & V & F & V & F & F! & F & V & V & F & F & V
\end{tabular}
\]

\end{solucio}

\begin{exercici}
Formalitzeu els següents enunciats: (1) \enquote{Tots els enters
que no són senars són parells}; (2) \enquote{Hi ha un nombre enter que no és parell}; (3)
\enquote{Per a cada nombre real $x$, hi ha un nombre real $y$ per al
qual $y^{3}=x$}; i (4) \enquote{Donats dos nombres
racionals $a$ i $b$, el producte $ab$ és racional}.
\end{exercici}

\begin{solucio}
Primer identifiquem els predicats i després escrivim l'enunciat simbòlicament.

(1) Suposem que $P$ és el predicat \enquote{ser parell} i $S$ és el predicat
"ser senar", aleshores es té:%
\[
\forall x\in\mathbb{Z}\text{, }\lnot S(x)\longrightarrow P(x)\text{,}%
\]
o també, més formalment,
\[
\left(  \forall x\right)  \left(  \left(  x\in\mathbb{Z}\wedge\lnot
S(x)\right)  \longrightarrow P\left(  x\right)  \right)  \text{.}%
\]


(2) Usant la mateixa notació que l'apartat anterior, es té:%
\[
\exists x\in\mathbb{Z}\text{, }\lnot P\left(  x\right)  \text{,}%
\]
o bé,%
\[
\left(  \exists x\right)  \left(  x\in\mathbb{Z\wedge}\lnot P\left(  x\right)
\right)  \text{.}%
\]


(3) L'enunciat simbòlicament s'escriu així:%
\[
\forall x\in\mathbb{R}\text{, }\exists y\in\mathbb{R}\text{, }y^{3}=x\text{,}%
\]
o bé, també com%
\[
\left(  \forall x\right)  \left(  \exists y\right)  \left(  x\in
\mathbb{R\wedge}\text{ }y\in\mathbb{R\wedge~}y^{3}=x\right)  \text{.}%
\]


(4) L'enunciat s'escriu així simbòlicament:%
\[
\forall a,b\in\mathbb{Q}\text{, }ab\in\mathbb{Q}\text{,}%
\]
o bé, com%
\[
\left(  \forall a,b\right)  \left(  a,b\in\mathbb{Q}\longrightarrow
ab\in\mathbb{Q}\right)  \text{.}%
\]

\end{solucio}

\begin{exercici}
Expresseu formalment els enunciats següents:

\begin{enumerate}
\item Algú no va fer els deures.

\item Tothom al dormitori té un company de pis que no li agrada.

\item Si algú al dormitori té un amic que té el xarampió,
llavors tothom el dormitori s'haurà de posar en quarantena.

\item Tot en aquesta botiga té un preu excessiu o està mal fet.
\end{enumerate}
\end{exercici}

\begin{solucio}
(1) Si $P(x)=$\enquote{$x$ fa els deures}, aleshores
l'enunciat es formalitza d'aquesta manera: $\exists x\lnot P(x)$.

(2) Necessitem tres predicats per formalitzar l'enunciat: $P(x)=$%
\enquote{$x$ viu en el pis}, $Q(x,y)=$%
\enquote{$x$ i $y$ són company de pis}\ i
$R(x,y)=$\enquote{$x$ li agrada $y$}. Aleshores es
té:%
\[
\forall x(P(x)\longrightarrow\exists y\left(  Q\left(  x,y\right)  \wedge\lnot
R\left(  x,y\right)  \right)
\]


(3) Introduïm quatre predicats per formalitzar l'enunciat: $P(x)=$%
\enquote{$x$ viu en el pis}, $Q(x,y)=$%
\enquote{$x$ i $y$ són amics}, $R(x)=$%
\enquote{$x$ té xarampió}\ i $S(x)=$%
\enquote{$x$ està en quarantena}. Aleshores
escrivim%
\[
\forall x\left(  P(x)\longrightarrow\exists y\left(  \left(  R(y)\wedge
Q\left(  x,y\right)  \right)  \longrightarrow S(x)\right)  \right)
\]


(4) Si considerem $P(x)=$\enquote{$x$ està a la
botiga}, $Q(x)=$\enquote{$x$ té el preu
excessiu}\ i $R(x)=$\enquote{$x$ està mal
fet}, llavors l'enunciat formalitzat s'escriu així:%
\[
\forall x\left(  P(x)\longrightarrow\left(  Q(x)\vee R(x)\right)  \right)
\]

\end{solucio}

\begin{exercici}
Tradueix les expressions següents en paraules:

\begin{enumerate}
\item $\forall x\left(  \left(  H(x)\wedge\lnot\exists yC(x,y)\right)
\longrightarrow\lnot F(x)\right)  $, si $H(x)=$\enquote{x és un
home}, $C(x,y)=$\enquote{$x$ està casat amb $y
$}\ i $F(x)=$\enquote{$x$ és
feliç}.

\item $\forall x\left(  W(x)\wedge Z(x)\right)  $, si $Z(x)=$\enquote{%
$x$ li agrada les croquetes}\ i $W(x)=$\enquote{$x$
té una granota de mascota}.

\item $\exists x\left(  N(x)\wedge\forall y\left(  N(y)\longrightarrow x\leq
y\right)  \right)  $, si $N(x)=$\enquote{$x$ és un nombre
natural}.
\end{enumerate}
\end{exercici}

\begin{solucio}
(1) \enquote{Tot home que no està casat no és
feliç}, o dit en unes altres paraules, \enquote{Tots els homes solters són infeliços}.

(2) \enquote{A tots els que tenen una granota com a mascota els
agraden les croquetes}.

(3) \enquote{Existeix un nombre natural més petit o igual que
qualsevol altre nombre natural}.
\end{solucio}

\begin{exercici}
Què signifiquen les afirmacions següents? Són vertaderes o falses?
L'univers del discurs en cada cas és $\mathbb{N}$, el conjunt de tots els
nombres naturals.

\begin{enumerate}
\item $\forall x\exists y(x<y)$.

\item $\exists y\forall x(x<y)$.

\item $\exists x\exists y(x<y)$.

\item $\exists x\forall y(x<y)$.

\item $\forall y\forall x(x<y)$.
\end{enumerate}
\end{exercici}

\begin{solucio}
(1) Això vol dir que per a cada nombre natural $x$, l'afirmació
$\exists y(x<y)$ és certa. En altres paraules, per a cada nombre natural
$x$, hi ha un nombre natural més gran que $x$. Això és cert. Per
exemple, $x+1$ sempre és més gran que $x$.

(2) Això vol dir que hi ha algun nombre natural $y$ tal que l'enunciat
$\forall x(x<y)$ és cert. En altres paraules, hi ha algun nombre natural
$y$ tal que tots els nombres naturals són més petits que $y$. Es clar
que això és fals. No importa el nombre natural $y$ que triem, sempre
hi haurà nombres naturals més grans.

(3) Això vol dir que hi ha un nombre natural $x$ tal que $\exists y(x<y)$
és cert. Però com hem vist al primer apartat, això és cert per
a tots els nombres naturals $x$, de manera que en particular és cert per
almenys un.

(4)\ Això vol dir que hi ha un nombre natural $x$ tal que l'enunciat
$\forall y(x<y)$ és cert. Podríeu tenir la temptació de dir que
aquesta afirmació serà certa si $x=0$, però això no és
correcte. Com que $0$ és el nombre natural més petit, l'enunciat $0<y
$ és cert per a tots els valors de $y$ excepte $y=0$, però si $y=0$,
llavors $0<y$ és fals i, per tant, $\forall y(0<y)$ és fals. Un
raonament similar mostra que per a cada valor de $x$ l'enunciat $\forall
y(x<y)$ és fals, per tant $\exists x\forall y(x<y)$ és fals.

(5) Això vol dir que per a cada nombre natural $x$, l'enunciat $\forall
y(x<y)$ és vertader. Però com hem vist a l'apartat anterior, ni tan
sols hi ha un valor de $x$ per al qual aquesta afirmació és certa.
Així doncs $\forall x\forall y(x<y)$ és fals.
\end{solucio}

\begin{exercici}
Escriu la negació de l'enunciat següent: \enquote{per a cada
nombre real $\varepsilon>0$, hi ha un enter positiu $k$ tal que per a tots els
enters positius $n$, es té $\left\vert a_{n}-k^{2}\right\vert
<\varepsilon$}.
\end{exercici}

\begin{solucio}
Considerem com univers del discurs el conjunt de tots els nombres reals i
simbolitzem per $\mathbb{Z}^{+}$ el conjunt del enters positius. D'aquesta
manera l'enunciat s'escriu formalment així:
\[
\forall\varepsilon>0,\exists k\in\mathbb{Z}^{+},\forall n\in\mathbb{Z}%
^{+},\left\vert a_{n}-k^{2}\right\vert <\varepsilon\text{.}%
\]


Aleshores la seva negació és:%
\[
\exists\varepsilon>0,\forall k\in\mathbb{Z}^{+},\exists n\in\mathbb{Z}%
^{+},\left\vert a_{n}-k^{2}\right\vert \geq\varepsilon\text{,}%
\]
que s'expressa en paraules com existeix un nombre real $\varepsilon>0$ tal que
per a tot nombre enter positiu $k$ existeix un nombre enter positiu $n$ que fa
que es compleixi $\left\vert a_{n}-k^{2}\right\vert \geq\varepsilon$.
\end{solucio}

\begin{exercici}
Assumint com a domini de les variables $x$ i $y$ el conjunt $U$, escriviu una
deducció per el següent raonament:%
\[%
\begin{tabular}
[c]{ll}%
$P_{1}:$ & $\forall x\left(  \left(  A(x)\longrightarrow R(x)\right)  \vee
T(x)\right)  $\\
$P_{2}:$ & $\exists x\left(  T(x)\longrightarrow P(x)\right)  $\\
$P_{3}:$ & $\forall x\left(  A(x)\wedge\lnot P(x)\right)  $\\\hline
$C:$ & $\exists x\ R(x)$%
\end{tabular}
\]

\end{exercici}

\begin{solucio}
Fem ara la deducció per provar la seva validesa:%
\[%
\begin{tabular}
[c]{lll}%
1. & $\forall x\left(  \left(  A(x)\longrightarrow R(x)\right)  \vee
T(x)\right)  $ & $P_{1}$\\
2. & $\exists x\left(  T(x)\longrightarrow P(x)\right)  $ & $P_{2}$\\
3. & $\forall x\left(  A(x)\wedge\lnot P(x)\right)  $ & $P_{3}$\\
4. & $T(a)\longrightarrow P(a)$ & EQE 2\\
5. & $\left(  A(a)\longrightarrow R(a)\right)  \vee T(a)$ & EQU 1\\
6. & $A(a)\wedge\lnot P(a)$ & EQU 3\\
7. & $\lnot P(a)$ & EC 6\\
8. & $\lnot T(a)$ & MT(4,7)\\
9. & $A(a)\longrightarrow R(a)$ & ED(5,8)\\
10. & $A(a)$ & EC 6\\
11. & $R(a)$ & MP(9,10)\\\hline
14. & $\exists x~R(x)$ & IQE 11
\end{tabular}
\]

\end{solucio}

\begin{exercici}
Escriviu una deducció del argument següent: \enquote{Tota
panerola intel·ligent menga escombraries. Hi ha una
panerola que li agrada la brutícia i no li agrada la pols. Per a cada
panerola, no és el cas que no li agradi la brutícia o mengi
escombraries. Per tant, hi ha una panerola tal que no és el cas que si no
és intel·ligent, li agradi la pols.}.
\end{exercici}

\begin{solucio}
Formalitzem els enunciats del raonament. Per facilitar l'escriptura,
considerem que la variable $x$ té per domini el conjunt de totes les
paneroles. Denotem per $I$, $E$, $B$ i $P$ els predicats \enquote{%
és intel·ligent}, \enquote{mengen porqueria}, \enquote{agrada la
brutícia}, i \enquote{agrada la
pols}, respectivament. Aleshores, el raonament podem
simbolitzar-lo d'aquesta manera:%
\[%
\begin{tabular}
[c]{ll}%
$P_{1}:$ & $\forall x~\left(  I(x)\longrightarrow E(x)\right)  $\\
$P_{2}:$ & $\exists x~\left(  B(x)\wedge\lnot P(x)\right)  $\\
$P_{3}:$ & $\forall x~\lnot\left(  \lnot B(x)\vee E(x)\right)  $\\\hline
$C:$ & $\exists x~\lnot\left(  \lnot I(x)\longrightarrow P(x)\right)  $%
\end{tabular}
\]
Fem ara la deducció per provar la seva validesa:%
\[%
\begin{tabular}
[c]{lll}%
1. & $\forall x~\left(  I(x)\longrightarrow E(x)\right)  $ & $P_{1}$\\
2. & $\exists x~\left(  B(x)\wedge\lnot P(x)\right)  $ & $P_{2}$\\
3. & $\forall x~\lnot\left(  \lnot B(x)\vee E(x)\right)  $ & $P_{3}$\\
4. & $B(a)\wedge\lnot P(a)$ & EQE 2\\
5. & $\lnot\left(  \lnot B(a)\vee E(a)\right)  $ & EQU 3\\
6. & $B(a)\wedge\lnot E(a)$ & DM 5\\
7. & $I(a)\longrightarrow E(a)$ & EQU 1\\
8. & $\lnot E(a)$ & EC 6\\
9. & $\lnot I(a)$ & MT(7,8)\\
10. & $\lnot P(a)$ & EC 4\\
11. & $\lnot I(a)\wedge\lnot P(a)$ & IC(9,10)\\
12. & $\lnot\left(  I(a)\vee P(a)\right)  $ & DM 11\\
13. & $\lnot\left(  \lnot I(a)\longrightarrow P(a)\right)  $ & EQ\\\hline
14. & $\exists x~\lnot\left(  \lnot I(x)\longrightarrow P(x)\right)  $ & IQE
13
\end{tabular}
\]
on EQ significa que hem aplicat l'equivalencia lògica: $A\longrightarrow
B\Longleftrightarrow\lnot A\vee B$.
\end{solucio}

\section{Demostració}

\begin{exercici}
Donat un enter positiu $n$, proveu directament que $n^{3}-n$ és sempre
múltiple de 3.
\end{exercici}

\begin{solucio}
Observa primer que $n(n^{2}-1)=n(n+1)(n-1)$. Aquest tres factors són tres
nombres naturals consecutius i, per tant, un d'ells serà un múltiple
de 3. Com a conseqüència $n^{3}-n$ és sempre múltiple de 3.
\end{solucio}


\begin{exercici}
Demostreu directament que per a cada terna de nombres reals positius $a,b$ i
$c$ es compleix que
\[
\frac{a}{b+c}+\frac{b}{a+c}+\frac{c}{a+b}\geq1.
\]
\end{exercici}

\begin{solucio}
De fet, provarem la desigualtat més forta
\[
\frac{a}{b+c}+\frac{b}{a+c}+\frac{c}{a+b}\geq \frac32.
\]
Com que $a,b,c>0$, podem escriure
\[
\frac{a}{b+c}=\frac{a^2}{a(b+c)},\qquad
\frac{b}{a+c}=\frac{b^2}{b(a+c)},\qquad
\frac{c}{a+b}=\frac{c^2}{c(a+b)}.
\]
Aplicant la desigualtat de Cauchy-Schwarz en la forma d'Engel, obtenim
\[
\frac{a^2}{a(b+c)}+\frac{b^2}{b(a+c)}+\frac{c^2}{c(a+b)}
\geq
\frac{(a+b+c)^2}{a(b+c)+b(a+c)+c(a+b)}.
\]
Ara bé,
\[
a(b+c)+b(a+c)+c(a+b)=2(ab+ac+bc).
\]
Per tant,
\[
\frac{a}{b+c}+\frac{b}{a+c}+\frac{c}{a+b}
\geq
\frac{(a+b+c)^2}{2(ab+ac+bc)}.
\]
Finalment, com que
\[
(a+b+c)^2=a^2+b^2+c^2+2(ab+ac+bc)\geq 3(ab+ac+bc),
\]
s'obté
\[
\frac{(a+b+c)^2}{2(ab+ac+bc)}\geq \frac32>1.
\]
Això demostra, en particular, la desigualtat demanada.
\end{solucio}

\begin{exercici}
Proveu cap a enrere que per a qualssevol nombres reals negatius, $a<b$ implica
$a^{2}>b^{2}$.
\end{exercici}

\begin{solucio}
La prova cap enrere surt de la tesi, i, per tant, podem fer es següent:%
\[%
\begin{tabular}
[c]{ll}%
$a^{2}>b^{2}$ & Tesi\\
$\left\vert a\right\vert >\left\vert b\right\vert $ & Prenent arrels
quadrades\\
$-a>-b$ & $a,b$ són negatius\\
$a<b$ & Multiplicant per $-1$ i surt l'hipòtesi
\end{tabular}
\]
De fet, això que hem escrit s'havia d'haver pensat mentalment perquè
la demostració real és una prova directa: Si $a<b$, multiplicant per
$-1$, es té $-a>-b$ i $-a,-b$ són positius. Per tant, $\left(
-a\right)  ^{2}>(-b)^{2}$ o sigui, $a^{2}>b^{2}$.
\end{solucio}

\begin{exercici}
Proveu per contrarecíproc que per a qualssevol $n,m\in\mathbb{Z}$, si
$mn$ és senar, llavors $m$ i $n$ són senars.
\end{exercici}

\begin{solucio}
Si fessim la prova directa seria prendre com hipòtesi que $mn$ és
senar i hem d'arribar a veure que $m$ i $n$ també ho són. En canvi,
per contrarecíproc serà prendre com hipòtesi que no és el cas
que $m$ i $n$ siguin senars i hem d'arribar a provar que $mn$ no és senar.

Si no és el cas que $m$ i $n$ siguin senars, vol dir que $m$ és parell
o $n$ és parell. Suposem per exemple que $m$ és parell. Aleshores
existeix $k\in\mathbb{Z}$ tal que $m=2k$. Per tant, $mn=2kn$ i $kn\in
\mathbb{Z}$. Aleshores $mn$ és parell com volíem demostrar.
\end{solucio}

\begin{exercici}
Proveu per reducció a l'absurd que els únics enters consecutius no
negatius $a,b$ i $c$ que compleixen $a^{2}+b^{2}=c^{2}$ són $3,4$ i $5$.
\end{exercici}

\begin{solucio}
Suposem que $3,4$ i~$5$ no són els únics enters consecutius no
negatius que compleixen $a^{2}+b^{2}=c^{2}$. Això és equivalent a
suposar que existeixen enters no negatius $n,n+1$ i $n+2$ i $n\neq3$ tals que%
\[
n^{2}+(n+1)^{2}=(n+2)^{2}%
\]
Ara bé, fent operacions, s'obté l'equació $n^{2}-2n-3=0$, què
té com a solucions $3$ i $-1$. Per tant, arribem a un absurd perquè
hem suposat que $n\neq3$. Com a conseqüència, hem provat el que volíem.
\end{solucio}

\begin{exercici}
Siguin $a,b$ i $c$ enters. Suposem que hi ha un nombre enter $d$ que $d\mid a
$ i $d\mid b$, però que $d$ no divideix $c$. Demostreu per reducció a
l'absurd que l'equació $ax+by=c$ no té cap solució que $x$ i $y $
siguin enters.
\end{exercici}

\begin{solucio}
Suposem que $ax+by=c$ té una solució tal que $x$ i $y$ són enters.
Com $d$ divideix $a$ i també $b$, aleshores existeixen enters $m$ i $n$
tals que $a=md$ i $b=nd$. Aleshores, es té%
\[
mdx+ndy=c
\]
Dividint per $d$, s'obté%
\[
mx+ny=\frac{c}{d}%
\]
Però això és absurd perquè $mx+ny$ és enter i, per tant,
també $\dfrac{c}{d}$ i això no és possible perquè $d$ no
divideix $c$.
\end{solucio}

\begin{exercici}
Proveu per inducció:

\begin{enumerate}
\item Si $n$ és un nombre enter no negatiu, llavors $%
%TCIMACRO{\dsum \limits_{k=0}^{n}}%
%BeginExpansion
{\displaystyle\sum\limits_{k=0}^{n}}
%EndExpansion
k\cdot k!=\left(  n+1\right)  !-1$.

\item Si $n\in\mathbb{N}$, llavors $\left(  1+x\right)  ^{n}\geq1+nx$ per a
tot $x\in\mathbb{R}$ i $x>-1$.
\end{enumerate}
\end{exercici}

\begin{solucio}
(1) Construïm una prova d'inducció sobre $n$, essent $P(n)=$
\enquote{$%
%TCIMACRO{\dsum \limits_{k=0}^{n}}%
%BeginExpansion
{\displaystyle\sum\limits_{k=0}^{n}}
%EndExpansion
k\cdot k!=\left(  n+1\right)  !-1$}.

\begin{description}
\item[Cas base:] $P(0)$ és $%
%TCIMACRO{\dsum \limits_{k=0}^{1}}%
%BeginExpansion
{\displaystyle\sum\limits_{k=0}^{0}}
%EndExpansion
k\cdot k!=\left(  0+1\right)  !-1$ i, això és veritat perquè
$0\cdot0!=0=1!-1$.

\item[Hipòtesi d'inducció:] Suposem que $P(n)$ és certa, o sigui
que $%
%TCIMACRO{\dsum \limits_{k=0}^{n}}%
%BeginExpansion
{\displaystyle\sum\limits_{k=0}^{n}}
%EndExpansion
k\cdot k!=\left(  n+1\right)  !-1$.

\item[Tesi:] Hem de demostrar que $P(n+1)=$\enquote{$%
%TCIMACRO{\dsum \limits_{k=0}^{n+1}}%
%BeginExpansion
{\displaystyle\sum\limits_{k=0}^{n+1}}
%EndExpansion
k\cdot k!=\left(  n+2\right)  !-1$}\ és certa.
\end{description}

En efecte, aplicant l'hipòtesi d'inducció, es té%
\begin{align*}%
%TCIMACRO{\dsum \limits_{k=0}^{n+1}}%
%BeginExpansion
{\displaystyle\sum\limits_{k=0}^{n+1}}
%EndExpansion
k\cdot k!  &  =%
%TCIMACRO{\dsum \limits_{k=0}^{n}}%
%BeginExpansion
{\displaystyle\sum\limits_{k=0}^{n}}
%EndExpansion
k\cdot k!+(n+1)(n+1)!\\
&  =\left(  n+1\right)  !-1+(n+1)(n+1)!\\
&  =(1+n+1)(n+1)!-1\\
&  =(n+2)!-1
\end{align*}
que és el que volíem veure.

\begin{description}
\item[Conclusió:] Pel principi d'inducció sobre $n$, deduïm que
$P(n)$ és certa per a tot $n\in\mathbb{N}$.
\end{description}

(2) Construïm una prova d'inducció sobre $n$, essent $P(n)=$
\enquote{Si $n\in\mathbb{N}$, llavors $\left(  1+x\right)  ^{n}%
\geq1+nx$ per a tot $x\in\mathbb{R}$ i $x>-1$}.

\begin{description}
\item[Cas base:] $P(1)$ és evident que és vertadera.

\item[Hipòtesi d'inducció:] Suposem que $P(n)$ és certa, o sigui
que si $n\in\mathbb{N}$, llavors $\left(  1+x\right)  ^{n}\geq1+nx$ per a tot
$x\in\mathbb{R}$ i $x>-1$.

\item[Tesi:] Hem de demostrar que $P(n+1)=$\enquote{Si $n\in
\mathbb{N}$, llavors $\left(  1+x\right)  ^{n+1}\geq1+(n+1)x$ per a tot
$x\in\mathbb{R}$ i $x>-1$}\ és certa.
\end{description}

En efecte, aplicant l'hipòtesi d'inducció, es té%
\begin{align*}
\left(  1+x\right)  ^{n+1}  &  =(1+x)^{n}(1+x)\\
&  \geq\left(  1+nx\right)  (1+x)\\
&  =1+x+nx+nx^{2}\\
&  =1+(n+1)x+nx^{2}\\
&  \geq1+(n+1)x
\end{align*}
que és el que volíem veure.

\begin{description}
\item[Conclusió:] Pel principi d'inducció sobre $n$, deduïm que
$P(n)$ és certa per a tot $n\in\mathbb{N}$.
\end{description}
\end{solucio}

\begin{exercici}
Proveu que el residu del quadrat de qualsevol nombre enter quan es divideix
per 4 és $0$ o $1$.
\end{exercici}

\begin{solucio}
Tot nombre enter $n$ és o parell o bé senar. Aquesta idea proporciona
l'estratègia de la demostració: construïr una prova per casos:

(1) Si $n$ és parell, aleshores $n=2k$, on $k\in\mathbb{Z}$. Llavors,
$n^{2}=4k^{2}$ què és múltiple de $4$ i, per tant, quan es
divideix per $4$ el residu val $0$.

(2) si $n$ és senar, aleshores $n=2k+1$, on $k\in\mathbb{Z}$. Llavors,
$n^{2}=4k^{2}+4k+1=4\left(  k^{2}+k\right)  +1$ què quan es divideix per
$4$ el residu val $1$.
\end{solucio}

\begin{exercici}
Demostreu que per a cada nombre real $x$, si $\left\vert x-3\right\vert >3$
llavors $x^{2}>6x$.
\end{exercici}

\begin{solucio}
Per definició de valor absolut, $\left\vert x-3\right\vert =x-3$ si
$x-3\geq0$, i $\left\vert x-3\right\vert =-\left(  x-3\right)  $ si $x-3<0$.
Construïm una prova per casos:

(1) Si $x-3\geq0$, aleshores $\left\vert x-3\right\vert =x-3>3$ i, per tant,
$x>6$. D'aquí, s'obté $x^{2}>6x$.

(2) Si $x-3<0$, aleshores $\left\vert x-3\right\vert =-x+3>3$ i, per tant,
$x<0$. D'aquí, s'obté $x^{2}>6x$.
\end{solucio}

\begin{exercici}
És cert que per a cada enter positiu $n$ es compleix que $n^{2}-n+17$
és un nombre primer?
\end{exercici}

\begin{solucio}
Si pensem que l'enunciat és fals hem de trobar un contraexemple per
provar-ho. En efecte, si prenem $n=17$, aleshores $n^{2}-n+17=289$ que no
és primer.
\end{solucio}


\begin{exercici}
(Algorisme de la divisió) Si $a,b\in\mathbb{N}$ i $b>0$, proveu que existeixen enters $q,r$ únics tals que $a=bq+r$, amb $0\leq r<b$.
\end{exercici}

\begin{solucio}
Considerem el conjunt
\[
S=\{a-bq:q\in\mathbb{Z},\ a-bq\geq 0\}.
\]
Aquest conjunt no és buit, perquè $a=a-b\cdot 0\in S$. Pel principi del mínim, $S$ té un element mínim; escrivim-lo com
\[
r=a-bq
\]
per a algun enter $q$. Per construcció, $r\geq 0$. Vegem que $r<b$. Si fos $r\geq b$, aleshores
\[
r-b=a-bq-b=a-b(q+1)\geq 0,
\]
i per tant $r-b\in S$. Però $r-b<r$, cosa que contradiu que $r$ era el mínim de $S$. Així, $0\leq r<b$ i $a=bq+r$.

Per provar la unicitat, suposem que
\[
a=bq+r=bq'+r',\qquad 0\leq r,r'<b.
\]
Aleshores
\[
r-r'=b(q'-q).
\]
Així, $r-r'$ és múltiple de $b$. Però, com que $0\leq r,r'<b$, tenim
\[
-b<r-r'<b.
\]
L'únic múltiple de $b$ que pot estar estrictament entre $-b$ i $b$ és $0$. Per tant, $r-r'=0$, és a dir, $r=r'$, i llavors també $q=q'$.
\end{solucio}
